\section{Uniform separation of the dyadic scales}
\label{sec:area_law_scale_separation}

Retain the fixed exponents and scales of
Definition~\ref{def:area_law_fine_belt_scales}. In particular,
$\delta_0=2\cdot10^{-7}$, $\zeta=1-\delta_0/2$, $t_k=2^{\ell_k}$,
$r_k=2^k$, and $s_k=2^{p_k}$. All indices below are nonnegative integers.
The following statements give the uniform scale separation used in
\cite[Section~11]{OpenAI2026AreaLaw}.

\begin{theorem}[Explicit gaps between scale indices]
    \label{thm:area_law_scale_index_gaps}
    \lean{TNLean.PEPS.AreaLaw.Geometry.dyadicScale_index_gaps}
    \leanok
    \uses{def:area_law_fine_belt_scales}
    For every pair of indices $d,k$ satisfying $10^7d\le k$,
    \begin{align}
        \ell_k+d & \le k,   \\
        k+d      & \le p_k.
    \end{align}
\end{theorem}
\begin{proof}\leanok
    \uses{thm:area_law_exponent_gaps}
    The threshold gives $d\le k$ and
    $\zeta k=k-10^{-7}k\le k-d$. Taking the floor proves the first
    inequality. Likewise $(1+\delta_0)k\ge k+2d\ge k+d$, so its floor
    is at least the integer $k+d$.
\end{proof}

\begin{theorem}[Dyadic factors between successive sides]
    \label{thm:area_law_scale_side_ratios}
    \lean{TNLean.PEPS.AreaLaw.Geometry.dyadicScale_side_ratios}
    \leanok
    \uses{def:area_law_fine_belt_scales}
    For every pair of indices $d,k$ satisfying $10^7d\le k$,
    \begin{align}
        2^d t_k & \le r_k, \\
        2^d r_k & \le s_k.
    \end{align}
\end{theorem}
\begin{proof}\leanok
    \uses{thm:area_law_scale_index_gaps}
    Apply the increasing function $n\mapsto2^n$ to the two index
    inequalities and use $2^{a+b}=2^a2^b$.
\end{proof}

\begin{theorem}[Uniform multiplicative separation]
    \label{thm:area_law_uniform_scale_separation}
    \lean{TNLean.PEPS.AreaLaw.Geometry.exists_dyadicScale_side_ratios}
    \leanok
    \uses{def:area_law_fine_belt_scales}
    For every real number $M$, there is a nonnegative integer $K$ such
    that, for every nonnegative integer $k\ge K$,
    \begin{align}
        M t_k & \le r_k, \\
        M r_k & \le s_k.
    \end{align}
    The threshold is chosen from $M$ and the fixed exponents alone,
    independently of the origin, domain, cut, and neighborhood radius.
\end{theorem}
\begin{proof}\leanok
    \uses{thm:area_law_scale_side_ratios}
    Choose a nonnegative integer $d$ with $M<2^d$, using the unboundedness
    of powers of two, and set $K=10^7d$. For every $k\ge K$, the preceding
    theorem gives $2^dt_k\le r_k$ and $2^dr_k\le s_k$. The cell sides are
    positive, so replacing $2^d$ by $M$ proves both inequalities.
\end{proof}
